\respcap{cap:agm}
\resp{agm-l1} Por indução, o conjunto $A$ das arestas aceitas está contido em alguma AGM. Ao aceitar $e=uv$, tome como $S$ a árvore atual de $u$ na floresta. Nenhuma aresta de $A$ atravessa o corte $(S,V\setminus S)$, e $e$ é a mais leve que o atravessa: uma aresta mais leve que o atravessasse teria sido examinada antes, quando os seus extremos já estavam em árvores diferentes (as árvores só crescem), e teria sido aceita. Pelo Lema~\ref{lema:corte}, $A\cup\{e\}$ está contido em alguma AGM. Ao final, a saída é acíclica por construção e conexa, pois uma aresta entre duas componentes finais teria sido aceita; logo é uma árvore geradora, e é mínima. (Também segue do Teorema~\ref{teo:gulosomat} aplicado à matroide gráfica.)

\resp{agm-l2} Invariante: as arestas de $T$ estão contidas em alguma AGM. Vale no início, com $T$ sem arestas. A cada passo, o corte $(V(T),V\setminus V(T))$ não é atravessado por nenhuma aresta de $T$, e o algoritmo de Prim escolhe a mais leve que o atravessa; pelo Lema~\ref{lema:corte}, o invariante se mantém. Depois de $n-1$ passos, $T$ é geradora e está contida numa AGM, logo é uma AGM.

\resp{agm-l3} Use um heap de mínimo com os vértices fora da árvore, com chave $chave(v)=$ peso da aresta mais leve de $v$ até a árvore, e um vetor de posições para reduzir chaves. São $n$ remoções do mínimo e, para cada aresta $uv$ examinada quando $u$ entra na árvore, no máximo uma redução de chave (uma operação \emph{subir}, em $O(\log n)$). O total é $O((n+m)\log n)=O(m\log n)$, pois um grafo conexo tem $m\ge n-1$.

\resp{agm-l4} Use o vetor $chave$, sem heap (pseudocódigo da §\ref{sec:prim}): em cada um dos $n$ passos, uma varredura de $O(n)$ encontra o vértice de menor chave fora da árvore, e depois as chaves dos vizinhos são atualizadas. (a) Com listas, as atualizações percorrem cada lista uma vez, $O(m)$ no total; o custo é $O(n^2+m)=O(n^2)$. (b) Com matriz, percorre-se a linha $u$, em $O(n)$, a cada passo: $O(n^2)$.

\resp{agm-l6} Sim, como no exercício anterior. E não é possível fazer melhor nessa representação: qualquer algoritmo correto precisa ler $\Omega(n^2)$ posições da matriz. Se alguma posição $(u,v)$ não fosse lida, um adversário poderia colocar ali uma aresta levíssima, que estaria em toda AGM, e o algoritmo erraria. Logo $O(n^2)$ é ótimo com matriz, e a versão com heap seria pior, $O(n^2\log n)$.

\resp{agm-l7} (i) Executar o Kruskal com as arestas em ordem \textbf{não crescente} de peso, o que é o guloso de matroides aplicado diretamente aos pesos $c$ (Teorema~\ref{teo:gulosomat}). (ii) Trocar os pesos por $c'=C-c$, com $C>\max c$ para mantê-los positivos, ou por $-c$, e executar qualquer algoritmo de AGM: toda árvore geradora tem exatamente $n-1$ arestas, então $c'(T)=(n-1)C-c(T)$, e minimizar $c'$ é maximizar $c$. (iii) Executar o Prim escolhendo sempre a aresta \emph{mais pesada} do corte. No grafo do Exemplo~\ref{ex:agm}, a árvore geradora máxima pesa 71.

\resp{agm-sim} Kruskal: aceita $AC\,(1)$, $DF\,(2)$, $BE\,(3)$ e $CF\,(4)$; rejeita $AD\,(5)$ e $CD\,(5)$, que fecham ciclo com $A,C,F,D$; aceita $BC\,(5)$; rejeita todas as arestas de peso 6. Prim a partir de $A$: $AC\,(1)$, $CF\,(4)$, $FD\,(2)$, $CB\,(5)$ e $BE\,(3)$. O custo é \textbf{15}. A AGM é única: cada aresta escolhida é a \emph{única} mais leve de algum corte --- $AC$ no corte $\{A\}$, $DF$ em $\{D\}$, $BE$ em $\{E\}$, $CF$ em $\{A,C\}$ e $BC$ em $\{B,E\}$ ---, e uma aresta assim pertence a toda AGM, pois uma AGM sem ela poderia ser melhorada pela troca do Lema~\ref{lema:corte}.

\resp{agm-unica} Suponha duas AGMs distintas, $T_1\ne T_2$, e seja $e$ a aresta de menor peso em $T_1\triangle T_2$; digamos que $e\in T_1$. $T_2+e$ tem um ciclo, que contém alguma aresta $f\in T_2\setminus T_1$, pois do contrário o ciclo estaria todo em $T_1$. Como $e$ é a menor aresta da diferença simétrica e os pesos são distintos, $c(e)<c(f)$; então $T_2+e-f$ é uma árvore geradora mais leve que $T_2$, uma contradição.

\resp{agm-ciclo} Suponha que $e=uv$ pertença a uma AGM $T$. Retirando $e$, $T$ se divide em dois lados; o ciclo de $G$ que contém $e$ precisa atravessar esse corte de novo, por outra aresta $f$, com $c(f)<c(e)$. Então $T-e+f$ é uma árvore geradora mais leve, uma contradição.

\resp{agm-gargalo} Suponha um caminho $P$ de $u$ a $v$ em $G$ cuja maior aresta pese menos que a maior aresta $e$ do caminho de $u$ a $v$ em $T$. Retirando $e$ de $T$, $u$ e $v$ ficam em lados diferentes; $P$ atravessa esse corte por alguma aresta $f$ com $c(f)<c(e)$. Então $T-e+f$ é mais leve que $T$, uma contradição.

\resp{agm-nova} Em $T+uv$ há um único ciclo, formado pelo caminho de $u$ a $v$ em $T$ e pela nova aresta. Encontre esse caminho com uma busca em $T$, em $O(n)$, tome a aresta mais pesada do ciclo e retire-a; se for a própria $uv$, a AGM não muda. Pela propriedade do ciclo, a aresta retirada não está em nenhuma AGM do novo grafo, e as outras arestas de $G$ fora de $T$ continuam excluídas pelos mesmos ciclos de antes.

\resp{agm-prim7} Prim a partir de $A$: $AD\,(5)$, $DF\,(6)$, $AB\,(7)$, $BE\,(7)$, $EC\,(5)$ e $EG\,(9)$. Kruskal: aceita $AD\,(5)$, $CE\,(5)$, $DF\,(6)$, $AB\,(7)$, $BE\,(7)$ e $EG\,(9)$, e rejeita $BC\,(8)$, $EF\,(8)$, $BD\,(9)$, $FG\,(11)$ e $DE\,(15)$. O custo da AGM é \textbf{39}.

\resp{agm-boruvka} Com pesos distintos, a aresta mais leve que sai de uma componente $C$ é a mais leve do corte $(C,V\setminus C)$, e pertence à AGM (que é única) pela propriedade do corte; logo todas as arestas escolhidas numa rodada pertencem à AGM e podem ser acrescentadas juntas, sem formar ciclo. Cada componente se junta a pelo menos uma outra em cada rodada, então o número de componentes pelo menos cai à metade, e há no máximo $\log_2n$ rodadas. Cada rodada examina todas as arestas uma vez e une as componentes, em $O(m)$ (ou $O(m\,\alpha(n))$ com Union-Find). Total $O(m\log n)$.

\resp{agm-iguais} Toda árvore geradora tem o mesmo peso, $(n-1)c$, e portanto é mínima. Basta encontrar uma árvore geradora qualquer, por exemplo a árvore de uma busca em largura ou em profundidade, em $O(n+m)$.

\resp{agm-12} Separe as arestas em duas listas, as de peso 1 e as de peso 2, em $O(m)$; essa é a ``ordenação''. Execute o Kruskal processando primeiro a lista de peso 1 e depois a de peso 2. O custo é o das operações de Union-Find, $O(m\,\alpha(n))$.

\resp{agm-leve} A aresta mais leve $e=uv$, única, é a mais leve do corte $(\{u\},V\setminus\{u\})$, e é a única mais leve desse corte; pela troca do Lema~\ref{lema:corte}, uma AGM sem ela poderia ser melhorada, então ela está em toda AGM. A segunda mais leve, $f=xy$, se única, também está em toda AGM (num grafo sem arestas paralelas): um dos extremos de $f$, digamos $y$, não é extremo da mais leve, e então $f$ é a única aresta mais leve do corte $(\{y\},V\setminus\{y\})$. A terceira, não necessariamente: ela pode fechar um triângulo com as duas primeiras.

\resp{agm-dijk} O triângulo com $AB\,(2)$, $BC\,(2)$ e $AC\,(3)$. A partir de $A$, a árvore de caminhos mínimos usa $AB$ e $AC$ (pois $\operatorname{dist}(A,C)=3<4$), com peso total 5; a AGM é $\{AB,BC\}$, com peso 4.

\resp{agm-parar} No melhor caso, as $n-1$ primeiras arestas da ordem formam uma árvore, e só elas são examinadas. No pior caso, a última aresta da ordem é necessária (por exemplo, é a única que liga um vértice ao resto), e todas as $m$ arestas são examinadas. A complexidade assintótica não muda, porque a ordenação, de custo $O(m\log m)$, é feita de qualquer forma.

\resp{agm-listas} São $2m$ operações \Find, cada uma em $O(1)$ (um ponteiro), e $n-1$ uniões. Com ``menor na maior'', cada vértice muda de representante no máximo $\log_2n$ vezes, pois o tamanho do seu conjunto pelo menos dobra a cada mudança; o custo total das uniões é, portanto, $O(n\log n)$. O laço custa $O(m+n\log n)$.

\resp{agm-floresta} O Kruskal, executado até o fim, devolve uma \emph{floresta geradora mínima}: uma AGM de cada componente. O Prim só alcança a componente da raiz; para obter a floresta, repita-o a partir de um vértice ainda não alcançado, até cobrir todos os vértices.

\resp{agm-segunda} Seja $T'$ a segunda melhor árvore e $e\in T'\setminus T$ a aresta de menor peso dessa diferença. Pelo argumento de troca, $T+e$ tem um ciclo com uma aresta $f\in T\setminus T'$, e $T+e-f$ não é mais pesada que $T'$; repetindo o argumento, conclui-se que alguma árvore da forma $T+e-f$ é uma segunda melhor. Algoritmo: para cada vértice $u$, faça uma busca em $T$ a partir de $u$ calculando, para todo $v$, a aresta de maior peso no caminho de $u$ a $v$ em $T$, $\max(u,v)$: $O(n^2)$ no total. Depois, para cada aresta $uv\notin T$, o custo da troca é $c(uv)-\max(u,v)$; a melhor troca dá a segunda melhor árvore, em $O(n^2+m)=O(n^2)$.

\resp{agm-sim3} Kruskal: aceita $AC\,(2)$, $CE\,(3)$, $AB\,(4)$ e $DE\,(4)$; rejeita $BC\,(5)$, que fecharia o ciclo $A,B,C$; aceita $EF\,(8)$; rejeita $BD\,(10)$ e $DF\,(11)$. Prim a partir de $F$: $FE\,(8)$, $EC\,(3)$, $CA\,(2)$, $ED\,(4)$ e $AB\,(4)$; quando $A$ entra, $B$ e $D$ empatam com chave 4, e entra $D$, cuja chave foi definida antes. O custo é \textbf{21}.

\resp{agm-corte} O primeiro corte é atravessado por $AC\,(2)$, $BC\,(5)$ e $BD\,(10)$; a mais leve, $AC$, está em alguma AGM (de fato, na do exercício anterior). O segundo é atravessado por $BD\,(10)$, $DE\,(4)$ e $DF\,(11)$, e a mais leve, $DE$, também está na AGM.
